\begin{houseviewbox}[结论：行业价格改善是必要条件，不是公司盈利的充分证据]
2025 年 6 月碳酸锂价格曾跌至约 5.8 万元/吨的低位；外部行业资料显示，2026 年 1 月、2 月与 7 月初的均价约为 15.60、14.96 与 16.01 万元/吨。价格环境较 2025 年低点已改善，且 2026H1 中国新能源车销量达到 744.6 万辆，需求方向支持锂产业链修复。但这些是行业条件：它们不能替代雅化自身的实现收入、销量、单位成本、库存结转或套保损益披露。
\end{houseviewbox}

\section{把行业信号拆成可传导、可验证的四个环节}

锂盐研究最常见的错误，是将现货价格变化直接乘以设计产能，得出公司利润。雅化的基准模型没有这么做。它从 2025A 的 69,494 吨已披露销量和 6.94 万元/吨收入代理出发，先设定 2026E 的销量、实现收入代理和毛利率，再以正式 H1 的经营披露反向验证。这样做的代价是模型较保守；好处是不会把没有披露的库存、原料锁价、资源自供和客户结构伪装成已实现利润。

\begin{exhibitbox}[表 10-1：行业指标到雅化 EPS 的传导纪律]
\begin{tabularx}{\exhibitboxwidth}{L{2.5cm}L{2.7cm}X X}
\toprule
环节 & 已知事实 & 本报告允许使用的方式 & 不允许直接推导的结论 \\
\midrule
商品环境 & 2026 年初至 7 月初行业价格约 15 万--16 万元/吨 & 作为情景方向与周期位置的外部校验 & 等同雅化实际 ASP 或确认期收入 \\
下游需求 & 2026H1 中国新能源车销量 744.6 万辆 & 支持需求仍在扩张的方向判断 & 推导公司订单、客户份额或出货量 \\
公司销量 & 2025A 锂产品 69,494 吨 & 作为 2026E 7.5/9.0/10.0 万吨的历史锚 & 将近 13 万吨设计口径当作当期销量 \\
单位经济 & 2025A 锂产品毛利率 12.22\% & 对比 2026E 毛利率假设的起点 & 将行业价格反弹全额转成毛利 \\
现金转化 & 2025A、2026Q1 OCF 均为负 & 对估值和概率施加折价 & 用利润预告替代现金流验证 \\
\bottomrule
\end{tabularx}
\end{exhibitbox}

\section{三个周期情景：价格方向相同，利润路径并不相同}

熊、基、牛情景的区别不只是“锂价高低”。基准情景假设销量从 2025A 的 6.95 万吨提高至 9.0 万吨，实现收入代理为 10.5 万元/吨，毛利率为 32\%；它没有以 16 万元/吨的行业均价代替公司实现价，也没有把未披露的资源成本优势列为额外收益。牛情景要求销量、成本与现金三项共同验证，因此不是简单的价格敏感性。

\begin{exhibitbox}[表 10-2：2026E 周期变量与合并利润的映射]
\begin{tabularx}{\exhibitboxwidth}{L{3.0cm}R{2.0cm}R{2.0cm}R{2.0cm}X}
\toprule
变量 & 熊 & 基准 & 牛 & 估值处理 \\
\midrule
锂盐销量 & 7.5 万吨 & 9.0 万吨 & 10.0 万吨 & 低于设计口径；不按满产 \\
实现收入代理 & 9.5 万元/吨 & 10.5 万元/吨 & 11.2 万元/吨 & House 假设，非公司 ASP \\
锂业务毛利率 & 25\% & 32\% & 35\% & 不额外加资源、库存或套保信用 \\
锂业务毛利 & 17.8 亿元 & 30.2 亿元 & 39.2 亿元 & 量×价×毛利率 \\
合并归母净利 & 15.0 亿元 & 20.1 亿元 & 24.7 亿元 & H1 预告锚加 H2 独立桥 \\
经营现金流 & -1.0 亿元 & 2.0 亿元 & 6.0 亿元 & 仅为情景，需 H1 验证 \\
\bottomrule
\end{tabularx}
\end{exhibitbox}

\section{需求、供给与库存：本报告的证据边界}

行业供需研究可解释价格为何反弹，却不能独立证明单家公司获得了多少份额或毛利。公司披露了李家沟权益与优先供应安排、KMC 设计产能以及锂盐设计能力，但未同时披露本期实际来料量、成本、自供比例、库存成本、客户条款和实际利用率。因此本报告对资源项目的处理是\textbf{经营期权而非估值资产}：基准目标价给零额外溢价，只有正式披露使其可以和销量、成本、现金流对应时，才允许它抬升牛情景概率。

\begin{sourcequalitybox}[正式 H1 对周期判断的最小验证包]
优先检查四项：第一，锂盐收入与销量或等效运营指标能否支持 H2 量价组合；第二，分部毛利及成本变动能否解释利润跃升；第三，应收、存货与经营现金流是否同步改善；第四，资源或项目披露能否给出实际而非设计层面的量、成本和期限。缺任一关键项时，行业景气只能作为背景，不能成为估值上修依据。
\end{sourcequalitybox}

\sourcenote{公司 2025 年报、2026Q1 报告与 2026H1 业绩预告；SMM、中国金属报公开行业资料及中汽协新能源车销量。完整来源登记见附录 B。}
